\addtocontents{toc}{\protect\Needspace{3\baselineskip}}
\originalchapter{辐射、测光与物质状态}\label{ch:radiation}
\originalsection{强度与可观测量}
天体通常不能送入实验室. 光的能量、频谱和方向分布便是推断温度、运动与组成的入口. 比强度 $I_\nu$ 定义为沿方向 $\boldsymbol n$ 的射线在时间 $\dd t$ 内, 穿过垂直射线的面积 $\dd A_\perp$, 在立体角 $\dd\Omega$ 和频段 $\dd\nu$ 中携带的能量除以这四个增量. 若使用固定面元 $\dd A$, 则 $\dd A_\perp=\cos\theta\dd A$. 因而通量为 $F_\nu=\int I_\nu\cos\theta\dd\Omega$. 各向同性强度在整球上的净通量为零, 在向外半球上的通量为 $\pi I_\nu$. 强度与通量的区别在恒星表面与观测者处尤其重要.

半径为 $R$ 的球对称光源具有光度 $L$, 真空中能量守恒给出 $F=L/(4\pi d^2)$. 没有吸收时同一条射线的 $I_\nu$ 不随距离改变: 物体所占立体角与接收到的通量都按 $d^{-2}$ 缩小. 宇宙膨胀另使 $I_\nu/\nu^3$ 沿自由传播保持不变, 积分后的面亮度按 $(1+z)^{-4}$ 衰减. 后一结论有两个来自光子能量和到达率的因子, 再有两个来自角距离和光度距离之比.

星等是对数亮度: $m_1-m_2=-2.5\log_{10}(F_1/F_2)$. 绝对星等 $M$ 定义在 $10\,\mathrm{pc}$, 故无消光且低红移时 $m-M=5\log_{10}(d/10\,\mathrm{pc})$. 波段星等不能直接等同于总光度; 用不同温度的恒星比较时须保留波段响应和总光度校正. 消光 $A_\lambda$ 使 $m_\lambda$ 增加, 色余 $E(B-V)=A_B-A_V$ 则刻画相对变红.

\originalsection{黑体与输运方程}
频率为 $\nu$ 的光子振子平均能量为 $h\nu/(\ee^{h\nu/\kb T}-1)$. 每单位体积、频率有 $8\pi\nu^2/c^3$ 个含两偏振的模, 故
\begin{equation}\label{eq:planck}
B_\nu(T)=\frac{2h\nu^3}{c^2}\frac1{\ee^{h\nu/\kb T}-1},\qquad
u_\nu=\frac{4\pi}{c}B_\nu.
\end{equation}
这个分布需要热平衡和光子数可改变, 于是光子化学势为零. $h\nu\ll\kb T$ 时得到 Rayleigh--Jeans 极限 $B_\nu\simeq2\nu^2\kb T/c^2$; 高频端由指数抑制控制. 由 $B_\lambda\dd\lambda=-B_\nu\dd\nu$ 可知两个峰值并不以 $\lambda\nu=c$ 对应, 因为谱密度的横轴不同.

沿射线路径 $s$ 的发射系数 $j_\nu$ 和吸收系数 $\alpha_\nu$ 给出 $\dd I_\nu/\dd s=-\alpha_\nu I_\nu+j_\nu$. 定义 $\dd\tau_\nu=\alpha_\nu\dd s$、源函数 $S_\nu=j_\nu/\alpha_\nu$, 解一阶线性方程得
\begin{equation}
I_\nu(\tau)=I_\nu(0)\ee^{-\tau}+\int_0^\tau S_\nu(t)\ee^{-(\tau-t)}\dd t.
\end{equation}
局域热平衡且真吸收主导时 $S_\nu=B_\nu$. 常源函数解为 $I_\nu=I_\nu(0)\ee^{-\tau}+B_\nu(1-\ee^{-\tau})$. 光学厚介质抹掉背景, 光学薄介质保留背景并加上线积分发射. 散射主导时源函数含辐射场自身, 不能直接以黑体代替.

\originalsection{电离平衡}
在稀薄、非简并、局域热平衡气体中, 中性原子、电离离子和电子的化学势满足 $\mu_0=\mu_1+\mu_e$. 非相对论粒子的配分因子含热德布罗意体积 $\lambda_T^3=(2\pi\hbar^2/m\kb T)^{3/2}$. 对重原子与离子近似等质量, 消去化学势可得
\begin{equation}\label{eq:saha}
\frac{n_en_1}{n_0}=\frac{2U_1}{U_0}
 \left(\frac{2\pi m_e\kb T}{h^2}\right)^{3/2}\ee^{-\chi/\kb T}.
\end{equation}
$U_i$ 是不含自由电子自旋的内部配分函数, $\chi$ 为电离能, 前面的 2 是电子自旋简并. 高温不是唯一的电离条件; 低密度也提高电离比例, 因为电离增加粒子可用的平动状态. 公式不能用于强辐照且碰撞不足的光电离气体, 这种情形须比较实际电离与复合速率.

\originalsection{电离关系的统计推导}
为了把 Saha 式与可测密度联系起来, 需说明能量零点. 对非简并粒子 $i$, 单粒子能量写为 $E_i=E_{i0}+p^2/(2m_i)$, Maxwell--Boltzmann 数密度为 $n_i=U_i\lambda_i^{-3}\exp[(\mu_i-E_{i0})/\kb T]$. 反应 $H\leftrightarrow p+e$ 的化学势守恒并不等于三个 $E_{i0}$ 之和为零. 采用中性基态零点, 电离产物的最低总能比它高 $\chi$, 因而消去化学势后有 $\exp(-\chi/\kb T)$. 重离子与原子的热波长比近 1, 留下电子平动因子及电子自旋权重, 得式\eqref{eq:saha}.

这一步也说明为什么没有把原子质量放到最终的巨大 $T^{3/2}$ 因子中. 电离额外创造一个可以在空间独立平动的轻电子, 可用状态数正是由它的热波长控制. 在给定密度下半电离路径不是 $\kb T=\chi$, 而是指数与电子平动状态数共同竞争. 对高密度气体, 电荷屏蔽、能级截断与简并会使这里的理想配分函数失效.

辐射扩散同样可以从频率输运明确构造平均值. 局域黑体且光学厚时 $F_\nu=-(4\pi/3\kappa_\nu\rho)(\partial B_\nu/\partial T)\nabla T$. 对频率积分后定义
\[
\frac1{\kappa_R}=\frac{\int_0^\infty\kappa_\nu^{-1}(\partial B_\nu/\partial T)\dd\nu}
{\int_0^\infty(\partial B_\nu/\partial T)\dd\nu}.
\]
低不透明度频段更容易运走能量, 因而是调和权重. 分母由 $\int B_\nu\dd\nu=\sigma T^4/\pi$ 求导, 恢复正文的扩散式. 频率平均不是单纯的数学便利, 它对应实际的能量输运方向.
